\chapter{Chapter06 — Fourier Degree Bridge}
%%% generated by the blueprint pipeline from FABL.Chapter06.F₂Polynomials.FourierDegreeBridge
\section{Overview}

Declarations of the Lean module \texttt{FABL.Chapter06.F₂Polynomials.FourierDegreeBridge}.

\section{Declarations}

\begin{theorem}[Fourier coefficient of a complemented function]
\lean{FABL.fourierCoeff_one_sub_div_two}
\proofsource{Analysis-of-Boolean-Functions-by-Ryan-ODonnell}{pdf-ffbe9c7b3306-p150-b001}
\label{FABL.fourierCoeff_one_sub_div_two}
\leanok
\uses{FABL.expect_monomial, FABL.fourierCoeff, FABL.monomial}
\difficulty{3}
\proofstep
  {FABL.Sign}
  {context}
  {Analysis-of-Boolean-Functions-by-Ryan-ODonnell}
  {pdf-ffbe9c7b3306-p014-b002,
    pdf-ffbe9c7b3306-p014-b008}
\proofstep
  {FABL.signValue}
  {context}
  {Analysis-of-Boolean-Functions-by-Ryan-ODonnell}
  {pdf-ffbe9c7b3306-p014-b002,
    pdf-ffbe9c7b3306-p014-b008}
\proofstep
  {FABL.monomial}
  {direct}
  {Analysis-of-Boolean-Functions-by-Ryan-ODonnell}
  {pdf-ffbe9c7b3306-p022-b001,
    pdf-ffbe9c7b3306-p022-b004}
\proofstep
  {FABL.signMonomialChar}
  {context}
  {Analysis-of-Boolean-Functions-by-Ryan-ODonnell}
  {pdf-ffbe9c7b3306-p023-b001,
    pdf-ffbe9c7b3306-p023-b004}
\proofstep
  {FABL.signMonomialChar_injective}
  {context}
  {Analysis-of-Boolean-Functions-by-Ryan-ODonnell}
  {pdf-ffbe9c7b3306-p023-b006}
\proofstep
  {FABL.fourierCoeff}
  {direct}
  {Analysis-of-Boolean-Functions-by-Ryan-ODonnell}
  {pdf-ffbe9c7b3306-p025-b004}
\proofstep
  {FABL.expect_monomial}
  {direct}
  {Analysis-of-Boolean-Functions-by-Ryan-ODonnell}
  {pdf-ffbe9c7b3306-p024-b007,
    pdf-ffbe9c7b3306-p025-b001}
Let $n$ be a nonnegative integer, let $f:\{-1,1\}^n\to\mathbb R$, and let $S$ be a
subset of the coordinates. With Fourier coefficients computed under the uniform measure
using the characters $\prod_{i\in S}x_i$, the coefficient of the function
$x\mapsto(1-f(x))/2$ at $S$ is $\bigl(\mathbf 1_{S=\varnothing}-\widehat
f(S)\bigr)/2$, where $\mathbf 1_{S=\varnothing}$ is $1$ when $S$ is empty and $0$
otherwise.
\end{theorem}

\begin{theorem}[Function algebraic degree boolean function f₂encoding le fourier degree]
\lean{FABL.functionAlgebraicDegree_booleanFunctionF₂Encoding_le_fourierDegree}
\leanok
\uses{FABL.BooleanFunction, FABL.BooleanFunction.toReal, FABL.booleanFunctionF₂Encoding, FABL.fourierDegree, FABL.functionAlgebraicDegree, FABL.functionAlgebraicDegree_booleanFunctionF₂Encoding_le_fourierDegree_viaPolynomial}
\statementsource{Analysis-of-Boolean-Functions-by-Ryan-ODonnell}{pdf-ffbe9c7b3306-p150-b002}
\difficulty{1}
% dataset-pipeline: stub -- prose pending build_informal_statements.py
Statement pending regeneration.
\end{theorem}
